\Needspace{8\baselineskip}\originalchapter{等离子体、屏蔽与尺度}\label{ch:plasma-basic}
\originalsection{电离与集体行为}
等离子体是含有可自由运动的带电粒子的物质体系. 它可以高度电离, 也可以只含少量自由电子和离子; 判断其行为不能只看电离百分比. 长程 Coulomb 力把大量粒子的运动耦合, 使无碰撞的等离子体也能支持集体波动. 中性气体的声波常靠短程碰撞传递, 电子等离子体振荡则靠电场恢复.

本册温度 T 默认指能量 $k_BT_K$, 常用 eV; $1\,\mathrm{eV}=1.602176634\times10^{-19}\,\mathrm J$, 等于约11604.5 K. 压强 $p=nT$ 时 T 必须是焦耳, 不能把 eV 数值直接代入 SI 方程. 多物种各有温度, $p=\sum_sn_sT_s$, 电子与离子并不总处于同温.

热力学平衡中电离可用 Saha 方程估算. 对单电离、适当简并因子版本,
\begin{equation}
 \frac{n_en_i}{n_0}=\frac{2(2\pi m_eT)^{3/2}}{h^3}\e^{-\chi/T}.
\end{equation}
右边量纲为数密度, 电离程度还依总密度与压强. 放电中的稳态常由碰撞电离、复合和壁损失平衡, 未必满足 Saha 的局部热平衡假设. ``温度低于13.6 eV''不能推出氢完全不电离, Maxwell 高能尾仍可越过阈值; 也不能据此断言所有密度下1 eV就充分电离.

\originalsection{Debye 屏蔽}
静电势中平衡物种满足 $n_s=n_{s0}\exp(-q_s\phi/T_s)$. 在 $|q_s\phi|\ll T_s$ 时线性展开, 背景电中性 $\sum_sq_sn_{s0}=0$, 代 Poisson 方程得
\begin{equation}\label{eq:debye}
 (\nabla^2-\lambda_D^{-2})\phi=-\rho_{\rm ext}/\epsilon_0,
 \qquad\lambda_D^{-2}=\sum_s\frac{n_{s0}q_s^2}{\epsilon_0T_s}.
\end{equation}
静止点电荷 Q 的解为 $\phi=Q\e^{-r/\lambda_D}/(4\pi\epsilon_0r)$. 电子积聚和离子排斥形成包围外电荷的屏蔽云. 若只允许电子移动, 则只保留电子项; ``离子冷''不等于``离子不动'', 因 Boltzmann 响应正比 $1/T_i$, 静态冷离子反而响应更强.

对半径 a 的导体球, 固定表面势 $\phi_s$, 外部线性解为 $\phi(r)=\phi_s(a/r)\e^{-(r-a)/\lambda_D}$. 电场法向积分给 Q=$4\pi\epsilon_0a\phi_s(1+a/\lambda_D)$, 等效电容比真空球大. 这依边界条件与线性响应, 并不是对强负偏压鞘层的通用电容公式.

\originalsection{准中性与等离子体参数}
准中性指宏观尺度上 $\sum_sq_sn_s$ 远小于各单独电荷密度, 不指所有地方严格零. 用 Poisson 标度, 若 $e\phi\sim T$ 且变化尺度 L, 则相对电荷不平衡约为 $(\lambda_D/L)^2$. 因而 L远大于屏蔽尺度时可以在主要动力学中使用准中性, 边界鞘层和高频电子波仍需 Poisson 方程.

Debye 球内粒子数 $N_D=(4\pi/3)n\lambda_D^3$. 本课有时把 $\Lambda=n\lambda_D^3$ 称为等离子体参数, 差一个明确几何因子. $N_D\gg1$ 表示平均场与小角累积碰撞描述可控. 耦合参数 $\Gamma=e^2/(4\pi\epsilon_0a_{
\rm WS}T)$, $a_{\rm WS}=(3/(4\pi n))^{1/3}$, $\Gamma\gtrsim1$ 时强耦合, 简单 Debye 和弱耦合碰撞近似失效.

电子量子简并还要看 $n\lambda_{\rm th}^3$, $\lambda_{\rm th}=h/\sqrt{2\pi m_eT}$. 高密度低温可需要量子统计, 这与强耦合不是同一个判据. 一个体系可能经典强耦合, 也可能量子弱耦合.

\originalsection{电子振荡}
把电子相对固定离子微移 $\xi$, 线性电场 $E=en\xi/\epsilon_0$ 产生 $m_e\ddot\xi=-eE$, 给 $\omega_{pe}^2=ne^2/(\epsilon_0m_e)$. 频率 $f_{pe}=\omega_{pe}/(2\pi)$, 不能把两者数值混用. 离子也可响应时, 对相对位移恢复频率会含电子和离子逆质量之和. 不同几何的表面场改变系数, 圆柱横移频率并不直接等于一维平板的 $\omega_{pe}$.
